\subsection[Graded bundles and $n$-fold graded bundles]{Graded bundles and $\boldsymbol{n}$-fold graded bundles}

An important class of graded manifolds are those that carry non-negative grading. For the moment we will consider only purely even manifolds explicitly, although the statements in this subsection generalise to the supercase. We will furthermore require that this grading is associated with a smooth action $h\colon \R\ti F\to F$ of the monoid $(\R,\cdot)$ of multiplicative reals on a mani\-fold~$F$; a \emph{homogeneity structure} in the terminology of~\cite{JG_MR_higher_vb}. This action reduced to $\R_{>0}$ is the one-parameter group of dif\/feomorphism integrating the \emph{weight vector field}, thus the weight vector f\/ield is in this case \emph{$h$-complete}~\cite{Grabowski2013} and only \emph{non-negative integer weights} are allowed. Thus the algebra $\mathcal{A}(F)\subset C^\infty(F)$ spanned by homogeneous functions is $\mathcal{A}(F) = \bigoplus_{i \in \mathbb{N}}\mathcal{A}^{i}(F)$, where $\mathcal{A}^{i}(F)$ consists of homogeneous function of degree~$i$.

Importantly, we have that for $t \neq 0$ the action $h_{t}$ is a dif\/feomorphism of~$F$ and, when $t=0$, it is a smooth surjection $\tau=h_0$ onto $F_{0}=M$, with the f\/ibres being dif\/feomorphic to~$\mathbb{R}^{N}$ (cf.~\cite{JG_MR_higher_vb}). Thus, the objects obtained are particular kinds of \emph{polynomial bundles} $\tau\colon F\to M$ (e.g.,~\cite{Bertram:2014, Voronov:2010_highr_alg}), i.e., f\/ibrations which locally look like $U\times\R^N$ and the change of coordinates (for a~certain choice of an atlas) are polynomial in $\R^N$. For this reason graded manifolds with non-negative weights \emph{and} $h$-complete weight vector f\/ields~$\zD$ are also known as \emph{graded bundles}~\cite{JG_MR_higher_vb}.

\begin{Example}\looseness=-1 If the weight is constrained to be either zero or one, then the weight vector f\/ield is precisely a vector bundle structure on~$F$ and will be generally referred to as an \emph{Euler vector field}.
\end{Example}
\begin{Example} The principle canonical example of a graded bundle is the higher tangent bundle~$\sT^{k}M$; i.e., the $k$-th jets (at zero) of curves $\gamma\colon \mathbb{R} \rightarrow M$. Given a~smooth function~$f$ on a~mani\-fold~$M$, one can construct functions $f^{(\alpha)}$ on $\sT^k M$, where $0\leq \alpha\leq k$, the so called \emph{$(\alpha)$-lifts} of $f$ (see~\cite{Morimoto_Lifts}). They are def\/ined by
\begin{gather*}
f^{(\alpha)}([\gamma]_k):= \left.\frac{\dd^\alpha}{\dd t^\alpha}\right|_{t=0} f(\gamma(t)).
\end{gather*}
where $[\gamma]_k\in \sT^k M$ is the class of the curve $\gamma\colon \R\rightarrow M$. The functions $f^{(k)}\colon \sT^k M\to \R$ and $f^{(1)}\colon \sT M\to \R$ are called the \emph{$k$-complete lift} and the \emph{tangent lift} of~$f$, respectively. A coordinate system~$(x^a)$ on $M$ gives rise to so called \emph{adapted coordinate systems} $\big(x^{a, (\alpha)}\big)_{0\leq \alpha \leq k}$ on $\sT^k M$ in which $x^{a, (\alpha)}$ is of degree~$\za$, i.e.,
\begin{gather*}\zD=\sum_{a,\za}\za x^{a, (\alpha)}\pa_{x^{a, (\alpha)}}.
\end{gather*}
\end{Example}


On a general graded bundle, one can always pick an atlas of $F$ consisting of charts for which we have homogeneous local coordinates $\big(x^{A},y_{w}^{i}\big)$, where $\w\!\big(x^{A}\big) =0$ and $\w(y_{w}^{i}) = w$ with $1\leq w\leq k$, for some $k \in \mathbb{N}$ known as the \emph{degree} of the graded bundle. Note that, according to this def\/inition, a graded bundle of degree $k$ is automatically a graded bundle of degree $l$ for $l\ge k$. However, there is always a \emph{minimal degree}.

It will be convenient to group all the coordinates with non-zero weight together, as these form a basis of the function algebra of the graded bundle. The index $i$ should be considered as a ``generalised index'' running over all the possible weights. The label $w$ in this respect largely redundant, but it will come in very useful when checking the validity of various expressions. The local changes of coordinates respect the weight and hence are polynomial for non-zero weight coordinates. A little more explicitly, the changes of local coordinates are of the form
\begin{gather*}
 x^{A'} = x^{A}(x), \qquad y^{i'}_{w} = \sum \frac{1}{n!} y^{j_{1}}_{w_{1}} y^{j_{2}}_{w_{2}} \cdots y^{j_{n}}_{w_{n}}T_{j_{n} \cdots j_{2} j_{1}}^{\:\:\:\:\:\:\:\:\:\:\:\: \:\:\:\: i'}(x), \end{gather*}
where $w=w_1+\cdots+w_n$ and we assume the tensors $T_{j_{n} \cdots j_{2} j_{1}}^{\:\:\:\:\:\:\:\:\:\:\:\: \:\:\:\: i'}$ to be symmetric in lower indices. Naturally, changes of coordinates are invertible and so, automatically, $\big(T_{j}^{\:\: i'}(x)\big)$ is an invertible matrix.

Importantly, a graded bundle of degree $k$ admits a sequence of surjections
\begin{gather}\label{eqn:fibrations}
F=F_{k} \stackrel{\tau^{k}_{k-1}}{\longrightarrow} F_{k-1} \stackrel{\tau^{k-1}_{k-2}}{\longrightarrow} \cdots \stackrel{\tau^{3}_2}{\longrightarrow} F_{2} \stackrel{\tau^{2}_1}{\longrightarrow}F_{1} \stackrel{\tau^{1}}{\longrightarrow} F_{0} = M,
\end{gather}
where $F_l$ itself is a graded bundle over $M$ of degree $l$, obtained from the atlas of $F_k$ by removing all coordinates of degree greater than~$l$ (see the next paragraph).

Note that $F_{1} \rightarrow M$ is a linear f\/ibration and the other f\/ibrations $F_{l} \rightarrow F_{l-1}$ are af\/f\/ine f\/ibrations in the sense that the changes of local coordinates for the f\/ibres are linear plus an additional additive terms of appropriate weight. The model f\/ibres here are~$\mathbb{R}^{n}$.

Morphisms between graded bundles necessarily preserve the weight; in other words, morphisms relate the respective homogeneity structures, or equivalently morphisms relate the respective weight vector f\/ields. Evidently, morphisms of graded bundles can be composed as standard maps between smooth manifolds and so we obtain the category of graded bundles. We will denote the category of graded bundles of degree $k$ by $\GrB[k]$.

\begin{Remark}
We will also encounter $N$-manifolds in this paper, which are very similar to graded bundles in many respects, and with hindsight are special examples of graded manifolds as def\/ined by Voronov \cite{Voronov:2001qf}. \v{S}evera \cite{Severa:2005} def\/ines an $N$-manifold as a supermanifold equipped with an action of $(\mathbb{R}, \cdot)$ such that $-1 \in \mathbb{R}$ acts as the parity operator (it f\/lips sign of any Grassmann odd coordinate). Roytenberg \cite{Roytenberg:2002} def\/ines an $N$-manifold in terms of an atlas of charts consisting of homogeneous coordinates for which coordinates of even weight are Grassmann even and coordinates of odd weight are Grassmann odd. It turns out that, under the additional assumption that even coordinates of non-zero degree are `cylindrical', both these notions of $N$-manifolds are equivalent, though the f\/irst (sketch) of a proof we know of is to be found in \cite{Bruce:2014c} and a complete proof can be found in \cite{Jozwikowski:2016}.
\end{Remark}

The notion of a double vector bundle \cite{Pradines:1974} (or an $n$-fold vector bundle \cite{Gracia-Saz:2009,Gracia-Saz:2012,Mackenzie:2005,Voronov:2012}) is conceptually clear in the graded language in terms of mutually commuting weight vector f\/ields; see \cite{Grabowski:2006,JG_MR_higher_vb}. This leads to the higher analogues known as \emph{$n$-fold graded bundles}, which are manifolds for which the structure sheaf carries an $\mathbb{N}^{n}$-grading such that all the weight vector f\/ields are $h$-complete and pairwise commuting. The local triviality of $n$-fold graded bundles means that we can always equip an $n$-fold graded bundle with an atlas such that the charts consist of coordinates that are simultaneously homogeneous with respect to the weights associated with each weight vector f\/ield (cf.~\cite{JG_MR_higher_vb}). If all the weight vector f\/ields are in fact Euler vector f\/ields, then we have an \emph{$n$-fold vector bundle}. The changes of local coordinates on an $n$-fold graded bundle must respect the weights. Similarly, morphisms between $n$-fold graded bundles respect the weights of the local coordinates.


We will denote a $k$-fold vector bundle as $\sD = (D; \Delta^1, \ldots, \Delta^k)$, where $\Delta^i$'s are Euler vector f\/ields on an underlying manifold $D$. We will often use the short hand `DVB' for double vector bundle. We shall use a similar notation for a $k$-fold graded bundle, but shall use the letter $\sF = (F; \Delta^1, \ldots, \Delta^k)$ to indicate that weights greater than one are allowed. We remark that the ordering of the weight vector f\/ields is a part of the def\/inition of a $k$-fold graded bundle and that morphisms of $k$-fold vector bundles respect this ordering. Thus, any permutation of the ordering of the weight vector f\/ields gives a dif\/ferent $k$-fold vector bundle structure on the same underlying manifold. The $k$-fold graded bundle $\sF$ gives rise to a number of graded bundles $\tau_i\colon F\rightarrow F_i$ def\/ined by pairs $(F; \Delta^i)$ called \emph{legs} of $\sF$, and a number of $(k-1)$-fold graded bundles $\sF_i = (F_i; \Delta^1, \ldots, \Delta^{i-1}, \Delta^{i+1}, \ldots, \Delta^k)$ which we call \emph{feet} of~$\sF$.

Important examples are: the \emph{iterated tangent bundle} $\sT^{(k)} M:=\sT\cdots\sT M$ (a $k$-fold vector bundle) and~$\sT^k\sT^l M$ (a double graded bundle).

By iterating the construction of $(\za)$-lifts, we obtain functions $f^{(\beta, \alpha)} := (f^{(\beta)})^{(\alpha)}$ on $\sT^k \sT^l M$ for $0\leq \alpha\leq k$, $0\leq \beta\leq l$, and, generally, functions $f^{(\ep_r, \ldots, \ep_1)}$ on $\sT^{n_1} \cdots \sT^{n_r} M$ for $0\leq \ep_j\leq n_j$, $1\leq j\leq r$. A coordinate system $(x^a)$ on $M$ gives rise to so called \emph{adapted coordinate sys\-tems}~$(x^{a, (\ep)})_\ep$ on $\sT^{(k)} M$, where the multi-index $\ep=(\ep_r, \ldots, \ep_1)$ is from $\{0,1\}^k$ and $x^{a, (\alpha)}$, $x^{a, (\ep)}$ are obtained from $x^a$ by the above lifting procedure.

\begin{Definition}\label{def:symmetricVB}
Consider a double vector bundle, which following classical notation we denote as $\sD=(D; A, B; M)$, where $A$ and $B$ are the side bundles. If we are given a vector bundle isomorphism $A \simeq B$ then $\sD$ is called \emph{balanced}. We may simply write $\sD=(D; A, A; M)$ but usually will avoid this. In more generality, a $k$-fold vector bundle~$\sD$ with legs $D\rightarrow D_i$, $i=1, \ldots, k$, is said to be a \emph{balanced $k$-fold vector bundle} if all feet $D_i$ are pairwise-isomorphic as $(k-1)$-fold vector bundles in some canonical way.
\end{Definition}

\subsection{Some constructions}
By saying that a vector f\/ield $\Delta$ on a manifold $F$ \emph{is a weight vector field} we mean that $F$ can be given a graded bundle structure which def\/ines $\Delta$ as its weight vector f\/ield. In particular, $\Delta$~has to be complete and the linear part of $\Delta$ in any of singular point of $\Delta$ has to have integer, non-negative eigenvalues. A important remark \cite[Remark~5.2]{JG_MR_higher_vb} is that the sum (unlike the dif\/ference) of two commuting weight vector f\/ields is again a weight vector f\/ield, hence a linear combination of weight vector f\/ields with non-negative integer coef\/f\/icients $a_s$, $\Delta = \sum_s a_s \Delta^s$ is a~weight vector f\/ield. We shall use this fact to set some further useful notation.

\begin{Definition}[removal coordinates of weight $> l$]\label{def:removal_coordinates1} Let $\sF = (F; \Delta)$ be a graded bundle.
Denote by $F[\Delta \leq l]$ the base manifold of the locally trivial f\/ibration def\/ined by taking the weight $>l$ coordinates with respect to this complementary weight to be the f\/ibre coordinates. We have a~natural projection that we will denote as
\begin{gather*}
\textnormal{p}^F_{[\Delta \leq l]}\colon \  F \rightarrow F[\Delta \leq l].
\end{gather*}
\end{Definition}
Note that, as the transformation rules of coordinates of weights $\le l$ involve only coordinates of weights $\le l$, the above def\/inition is correct. Since compositions of commuting homogeneity structures are homogeneity structures, we immediately get the following.
\begin{Proposition} Linear combinations of commuting weight vector fields with non-negative integer coefficients are weight vector fields. If at start $\sF = (F; \Delta^1, \ldots, \Delta^k)$ is a $k$-fold graded bundle and $\Delta=\sum a_s \Delta^s$, where $a_i\in\N$, then $F[\Delta \leq l]$ has the induced $k$-fold graded bundle structure such that $\textnormal{p}^F_{[\Delta \leq l]}$ is a $k$-fold graded bundle morphism.
\end{Proposition}

The next construction is in a sense dual and depends on putting to zero coordinates of negative weights with respect to the vector f\/ield $X=\sum a_s \Delta^s$. Here, we assume that $a_s\in \Zet$ can be negative (see Remark \ref{rem:non_integer_weights}).
\begin{Definition}[putting to zero coordinates of $X$-negative weights]\label{def:removal_coordinates2} Consider a $k$-fold graded bundle $(F; \Delta^1, \ldots, \Delta^k)$ and let $X=\sum_s a_s \Delta^s$ be a linear combination of weight vector f\/ields of~$F$ with integer coef\/f\/icients. Def\/ine locally a submanifold $F[X\geq 0] \subset F$ as
\begin{gather*}
 F[X\geq 0]:= \bigg\{\big(y^i_w\big)\colon y^i_w=0 \  \text{for all} \  w=(w_1, \ldots, w_k)\  \text{such that}\  \sum_s a_s w_s < 0\bigg\}.
\end{gather*}
\end{Definition}

Note f\/irst that $F[X\geq 0]$ is a well-def\/ined submanifold of $F$. Indeed, set $X$-weight $\w_X(f)$ for the weight with respect to the vector f\/ield $X$ of a homogeneous function $f$, thus $\w_X(y^i_w) = \sum_s a_s w_s$ if $w=(w_1, \ldots, w_s)$, and consider another coordinate system $(y^{i'})$. Since coordinate changes preserve $\w_X$ and $y^{i'}$ is a linear combination of monomials of type $y^{i_1}\cdots y^{i_j}$, we see that if $\w_X(y^{i'})$ is negative then at least one of the weight $\w_X(y^{i_1})$, $\w_X(y^{i_2}), \ldots, \w_X(y^{i_j})$ has to be negative, and thus $y^{i'}$ vanishes on $F[X\geq 0]$, hence the latter is invariantly def\/ined. As the weight vector f\/ields $\Delta^1, \ldots, \Delta^k$ preserve the $X$-weight, they are tangent to the submanifold $F[X\geq 0]$ which therefore has the induced $k$-fold graded bundle structure.

Considering $-X$ instead of $X$, we obtain graded subbundles $F[-X\ge 0]$, thus we can put to zero invariantly all coordinates of positive $X$-weight, and, eventually, graded subbundle \begin{gather*}
F[X=0] = F[X\geq 0] \cap F[-X\geq 0].
\end{gather*}
Is is clear that $F[X=0]$ consist of singular points of the vector f\/ield $X = \sum a_s \Delta^s$.

\begin{Remark}\label{rem:non_integer_weights} Formally we could admit even irrational coef\/f\/icient in $X$ to def\/ine $F[X\geq 0]$ or $F[X=0]$ but this does not enlarge possible constructions, e.g., if $\Delta^1 = x \pa_x$, and $\Delta^2 = y\pa_y$ are weight vector f\/ields of a double graded bundle $F$ then $F[x\pa_x + \sqrt{2} y \pa_y\geq 0]$ is the same as $F[x\pa_x + q y \pa_y\geq 0]$ for some rational number $q$ close to $\sqrt{2}$ since there is only a f\/inite number of weights assigned to coordinates; moreover, $F[X\geq 0] = F[k X\geq 0]$, $k\neq 0$, so we can forget about non-integer coef\/f\/icients in $X$.
\end{Remark}

\begin{Example}
If $(F, \Delta)$ is a graded bundle, then $F[\Delta=0]$ is the base of $F$.
\end{Example}
\begin{Example} \label{exmpl:core} If $\sD = (D; \Delta^1, \Delta^2)= (D; A, B; M)$ is a double vector bundle, then $C:=\sD[\Delta^1-\Delta^2=0]$ is the core bundle of $\sD$. Besides, $\sD[\Delta^1-\Delta^2\geq 0]$ is the kernel of the vector bundle morphism $D\rightarrow B$, so the intersection $\sD[\Delta^1-\Delta^2\geq 0]\cap \sD[\Delta^2-\Delta^1\geq 0]$ coincides with the core bundle.
\end{Example}

\begin{Remark} The projection $\textnormal{p}^F_{[\Delta \leq l]}$ from Def\/inition \ref{def:removal_coordinates1} and the inclusions
$F[X=0] \rightarrow F$ and $F[X\geq 0] \rightarrow F$ are functorial. We shall use this fact later on.
\end{Remark}

